% Informe de Animal Identification con Custom Vision, estilo editorial rayos-x.
\documentclass[a4paper,11pt]{article}

\usepackage[spanish]{babel}
\usepackage[T1]{fontenc}
\usepackage{newpxtext,newpxmath}
\usepackage{microtype}
\usepackage[skip=0.6\baselineskip]{parskip}
\usepackage[a4paper,margin=2.5cm]{geometry}
\usepackage{amsmath}
\usepackage{booktabs}
\usepackage{float}
\usepackage{graphicx}
\usepackage{url}
\usepackage{xcolor}
\usepackage{titlesec}
\usepackage{listings}

\definecolor{editorialink}{HTML}{292927}
\definecolor{editorialaccent}{HTML}{984C3D}
\usepackage[colorlinks=true,linkcolor=editorialaccent,citecolor=editorialaccent,
  urlcolor=editorialaccent,breaklinks=true,pdfauthor={Jheison Martinez Bolivar},
  pdftitle={Animal Identification con Custom Vision}]{hyperref}

\titleformat{\section}{\normalfont\Large\bfseries\color{editorialink}}
{\textcolor{editorialaccent}{\thesection}}{0.75em}{}
\titlespacing*{\section}{0pt}{1.45\baselineskip}{0.55\baselineskip}
\titleformat{\subsection}{\normalfont\large\bfseries\color{editorialink}}
{\thesubsection}{0.7em}{}
\titlespacing*{\subsection}{0pt}{1.15\baselineskip}{0.4\baselineskip}
\setlength{\emergencystretch}{2em}

% Listados de código de los anexos.
\lstset{basicstyle=\ttfamily\footnotesize,numbers=left,numberstyle=\tiny\color{gray},breaklines=true,frame=single,columns=fullflexible,keepspaces=true}

% Captura de pantalla centrada y compacta.
\newcommand{\figura}[2]{%
  \begin{figure}[H]
    \centering
    \includegraphics[width=0.75\linewidth]{capturas/#1}
    \caption{#2}
  \end{figure}
}

% Dos capturas lado a lado, para ahorrar espacio.
\newcommand{\figuraPar}[4]{%
  \begin{figure}[H]
    \centering
    \begin{minipage}[t]{0.48\linewidth}
      \centering
      \includegraphics[width=\linewidth]{capturas/#1}
      \caption{#2}
    \end{minipage}\hfill
    \begin{minipage}[t]{0.48\linewidth}
      \centering
      \includegraphics[width=\linewidth]{capturas/#3}
      \caption{#4}
    \end{minipage}
  \end{figure}
}

\title{Animal Identification con Custom Vision}
\author{Jheison Martinez Bolivar}
\date{\today}

\begin{document}

  \maketitle

  \begin{abstract}
    Este informe describe la creaci\'on de un clasificador de im\'agenes de
    elefantes, jirafas y leones con Azure Custom Vision~\cite{microsoft_customvision2024}
    sobre Azure AI Services~\cite{microsoft_cognitiveservices2024}. La
    suscripci\'on Azure for Students impone la policy
    \texttt{sys.regionrestriction}, que solo permite \texttt{francecentral},
    \texttt{belgiumcentral}, \texttt{northcentralus}, \texttt{spaincentral} y
    \texttt{chilecentral}; por eso el instructivo original en \texttt{East US}
    falla con \texttt{InvalidTemplateDeployment / RequestDisallowedByAzure}.
    Verificado por CLI, solo \texttt{northcentralus} soporta
    \texttt{CustomVision.Training} y \texttt{Prediction}, as\'i que los recursos
    \texttt{labunisalle} y \texttt{labunisalle-Prediction} se crearon all\'i en
    SKU S0. El proyecto \textit{Animal Identification} (Classification,
    Multiclass, dominio General [A2]) se entren\'o con 17 im\'agenes por clase y
    Quick Training: Precision, Recall y AP en 100\,\% sobre entrenamiento. Con 6
    im\'agenes nuevas de Wikimedia Commons el modelo publicado acierta 6 de 6
    con confianza media de 96.7\,\%.
  \end{abstract}

  %=========================================================================
  \section{Introducci\'on}
  \subsection{Objetivo}
  Entrenar un clasificador de im\'agenes en Custom Vision con fotos de
  elefantes, jirafas y leones, documentar el proceso (recursos, proyecto,
  etiquetas, entrenamiento, m\'etricas y publicaci\'on) y probar el modelo
  publicado con im\'agenes no vistas.

  \subsection{Restricci\'on regional}
  El instructivo pide la regi\'on \texttt{East US}, pero la suscripci\'on la
  bloquea. La policy \texttt{sys.regionrestriction} de la suscripci\'on solo
  permite cinco regiones, y de esas cinco solo \texttt{northcentralus} devuelve
  SKUs para \texttt{CustomVision.Training} y \texttt{Prediction}:

  \begin{lstlisting}
$ az policy assignment show --name sys.regionrestriction
Allowed: francecentral, belgiumcentral, northcentralus,
  spaincentral, chilecentral
$ az cognitiveservices account list-skus \
  --kind CustomVision.Training --location northcentralus
F0, S0
$ az cognitiveservices account list-skus \
  --kind CustomVision.Training --location francecentral
(vacio)
  \end{lstlisting}

  Por eso todo el trabajo se hace en \texttt{northcentralus}. El grupo de
  recursos \texttt{lab} puede seguir en \texttt{eastus}: la policy mira la
  ubicaci\'on del recurso, no la del grupo.

  %=========================================================================
  \section{Metodolog\'ia}
  \subsection{Recursos en Azure}
  Se crearon en el grupo \texttt{lab} dos recursos en SKU S0, verificados con
  \texttt{az cognitiveservices account list --resource-group lab} (salida en
  \texttt{capturas/01-cli-recursos.txt}):

  \begin{table}[H]
    \centering
    \caption{Recursos verificados por CLI.}
    \label{tab:recursos}
    \begin{tabular}{lll}
      \toprule
      Nombre & Kind & Location / SKU \\
      \midrule
      labunisalle & CustomVision.Training & northcentralus / S0 \\
      labunisalle-Prediction & CustomVision.Prediction & northcentralus / S0 \\
      \bottomrule
    \end{tabular}
  \end{table}

  \figura{02-portal-recursos.png}{Creaci\'on de los recursos en el portal de Azure.}

  \subsection{Proyecto de Custom Vision}
  El recurso no es el proyecto. El proyecto se crea en
  \texttt{customvision.ai} con el recurso de entrenamiento como contenedor de
  llaves. Ajustes usados:

  \begin{itemize}
    \item Name: \textit{Animal Identification}.
    \item Description: \textit{Image classification for animals}.
    \item Resource: \texttt{labunisalle [northcentralus]}.
    \item Project Type: Classification.
    \item Classification Type: Multiclass (un solo tag por imagen).
    \item Domain: General [A2].
  \end{itemize}

  \figura{03-create-project.png}{Creaci\'on del proyecto Animal Identification en customvision.ai.}

  \subsection{Carga y etiquetado}
  Se subi\'o por etiquetas \texttt{elephant}, \texttt{giraffe} y \texttt{lion},
  con 17 im\'agenes por clase. Cada lote se etiqueta al subirlo; una imagen con
  m\'as de una etiqueta romper\'ia el supuesto Multiclass, por eso se verifica
  un solo tag por imagen.

  \figura{04-tags.png}{Lotes etiquetados por clase, 17 im\'agenes cada uno.}

  \subsection{Entrenamiento}
  Se entren\'o con \texttt{Train} en modo Quick Training. La Iteraci\'on 1
  termin\'o el 8/11/2026 a las 6:51:53 PM con dominio General [A2], ID
  \texttt{71779344-fd60-4c43-9e8f-788ef5230282}.

  \figura{05-train.png}{Entrenamiento r\'apido lanzado y completado.}

  \figura{06-train-result.png}{Performance de la Iteraci\'on 1: Precision, Recall y AP globales en 100.0\,\%.}

  %=========================================================================
  \section{Resultados}
  \subsection{M\'etricas de entrenamiento}
  La Tabla~\ref{tab:performance} resume el panel Performance con umbral 50\,\%.
  El 100\,\% en las tres m\'etricas y en los tres tags se lee con cautela: con
  51 im\'agenes el modelo memoriza el set de entrenamiento; la prueba que
  importa es con im\'agenes nuevas.

  \begin{table}[H]
    \centering
    \caption{Performance por tag, Iteraci\'on 1.}
    \label{tab:performance}
    \begin{tabular}{lccc}
      \toprule
      Tag & Precision & Recall & Im\'agenes \\
      \midrule
      lion & 100.0\,\% & 100.0\,\% & 17 \\
      giraffe & 100.0\,\% & 100.0\,\% & 17 \\
      elephant & 100.0\,\% & 100.0\,\% & 17 \\
      \midrule
      Global & 100.0\,\% & 100.0\,\% & 51 \\
      \bottomrule
    \end{tabular}
  \end{table}

  \subsection{Prueba con im\'agenes nuevas}
  Se probaron 6 im\'agenes externas de Wikimedia Commons (carpeta
  \texttt{validacion/}, con fuentes y licencias en \texttt{FUENTES.txt}), no
  vistas en entrenamiento, con \textit{Quick Test} sobre la Iteraci\'on 1. Las
  6 clasifican bien, con confianza media de 96.7\,\%. El caso m\'as d\'ebil es
  la imagen con tres jirafas juntas (87.0\,\% con 11.4\,\% a elephant), la
  \'unica con oclusi\'on y m\'ultiples instancias.

  \begin{table}[H]
    \centering
    \caption{Quick Test con im\'agenes nuevas (6 de 6 aciertos).}
    \label{tab:quicktest}
    \begin{tabular}{lcc}
      \toprule
      Imagen & Predicci\'on & Confianza \\
      \midrule
      elephant\_val\_01 & elephant 98.0\,\% & 98.0\,\% \\
      elephant\_val\_02 & elephant 97.6\,\% & 97.6\,\% \\
      giraffe\_val\_01 & giraffe 98.5\,\% & 98.5\,\% \\
      giraffe\_val\_02 & giraffe 87.0\,\% & 87.0\,\% \\
      lion\_val\_01 & lion 99.8\,\% & 99.8\,\% \\
      lion\_val\_02 & lion 99.6\,\% & 99.6\,\% \\
      \bottomrule
    \end{tabular}
  \end{table}

  \figuraPar{07-quick-test-elephant1.png}{Elefante nuevo: 98.0\,\%.}{08-quick-test-elephant2.png}{Elefante nuevo: 97.6\,\%.}

  \figuraPar{09-quick-test-giraffe1.png}{Jirafa nueva: 98.5\,\%.}{10-quick-test-giraffe2.png}{Tres jirafas con oclusi\'on: 87.0\,\%.}

  \figuraPar{11-quick-test-lion1.png}{Le\'on nuevo: 99.8\,\%.}{12-quick-test-lion2.png}{Le\'on rugiendo: 99.6\,\%.}

  \subsection{Publicaci\'on y cliente}
  La Iteraci\'on 1 se public\'o como modelo \texttt{AnimalIdentification-v1}
  contra el recurso \texttt{labunisalle-Prediction} en \texttt{northcentralus}.
  El endpoint de imagen es:

  \begin{lstlisting}
https://labunisalle-prediction.cognitiveservices.azure.com/customvision/v3.0/Prediction/8745a5dd-92de-4fd7-967f-642ce3967c6c/classify/iterations/Iteration1/image
  \end{lstlisting}

  con headers \texttt{Prediction-Key: <clave del portal>} y
  \texttt{Content-Type: application/octet-stream}. La clave completa queda en
  el portal (\texttt{labunisalle-Prediction > Keys and Endpoint}) y no se
  imprime en este PDF por seguridad; se rota despu\'es de la entrega. La
  evidencia del endpoint queda en \texttt{capturas/13-publish-endpoint.txt}.

  La prueba del cliente se hizo con \texttt{curl} sobre dos im\'agenes de
  validaci\'on (salida completa en
  \texttt{capturas/14-cloudshell-test.txt}):

  \begin{lstlisting}
$ curl -s -X POST "$PREDICT_URL_IMAGE" \
  -H "Prediction-Key: $PKEY" \
  -H "Content-Type: application/octet-stream" \
  --data-binary "@validacion/lion/lion_val_01.jpg"
lion 0.9986, elephant 0.0008, giraffe 0.0005
$ curl -s -X POST "$PREDICT_URL_IMAGE" \
  --data-binary "@validacion/giraffe/giraffe_val_02.jpg"
giraffe 0.8700, elephant 0.1147, lion 0.0153
  \end{lstlisting}

  Los valores coinciden con el Quick Test del portal (99.8\,\% y 87.0\,\%),
  as\'i que el modelo publicado responde igual que el evaluado.

  %=========================================================================
  \section{Conclusiones}
  \begin{itemize}
    \item El bloqueo en \texttt{East US} no es un error del procedimiento sino
      de la suscripci\'on: la policy solo deja cinco regiones y solo
      \texttt{northcentralus} sirve para Custom Vision.
    \item Con 51 im\'agenes el 100\,\% en entrenamiento confirma
      memorizaci\'on del set; la evidencia de generalizaci\'on es el 6 de 6 en
      im\'agenes nuevas con media de 96.7\,\%.
    \item El \'unico caso d\'ebil combina oclusi\'on y m\'ultiples instancias
      (tres jirafas, 87.0\,\%), que es el l\'imite esperable de un
      clasificador de imagen completa sin detecci\'on.
  \end{itemize}

  \subsection{Limitaciones}
  El conjunto es peque\~no y no generaliza m\'as all\'a de este ejercicio. Las
  im\'agenes de validaci\'on vienen de Wikimedia Commons y se citan en
  \texttt{validacion/FUENTES.txt}. El paso siguiente ser\'ia una prueba m\'as
  grande y con fondos m\'as variados.

  \bibliographystyle{plain}
  \bibliography{bib/references}

  \appendix
  \section{Anexo: script de prueba del modelo}
  Listado completo de \texttt{anexos/probar\_prediccion.sh}, que prueba el
  modelo publicado con \texttt{curl}. La clave se exporta fuera del script.

  \lstinputlisting[language=bash]{anexos/probar_prediccion.sh}

\end{document}
